\chapter{Le travail du responsable de desk}\label{ch:fm:the-head-of-desks-job}

La réunion budgétaire de décembre fixe l'objectif de revenus de l'année suivante au résultat de cette année majoré de dix pour cent, et laisse les limites de risque du desk
telles qu'elles étaient. Le résultat de cette année est exactement ce que le desk attend de gagner~; ses revenus ont un ratio de Sharpe de 1.3. Avec ces limites, le desk
atteint le nouvel objectif dans 45 années sur cent. Le responsable de desk s'est vu confier un plan qui échoue plus souvent qu'il ne réussit, avant même la première transaction~;
pour l'atteindre trois années sur quatre, le desk devrait prendre plus de deux fois son risque. Ce chapitre porte sur les chiffres qu'un responsable de desk accepte en décembre
et suit tout au long de l'année.

\section{Le budget~: revenus, coûts et plan}

\begin{definition}[Budget de revenus, coût alloué]\label{def:fm:the-head-of-desks-job:budget}
Le \emph{budget de revenus}\index{budget de revenus} d'un desk est le revenu qu'il s'engage à produire l'année suivante, convenu avec la direction de la société,
et utilisé pour planifier ses coûts, son enveloppe de rémunération et son capital. Un \emph{coût alloué}\index{coût alloué} est une part des coûts communs de la société,
technologie, locaux, fonctions de contrôle, données de marché, imputée au desk selon une clé d'allocation (effectifs, revenus, usage).
\end{definition}

Le plan d'un desk comporte trois volets~: les revenus par activité, les coûts (ses coûts directs et ceux qui lui sont alloués), et le capital et les limites qu'il utilisera (chapitre~7).
Le desk du chapitre (illustratif, en millions de dollars par an) couvre trois activités~: le flux client, de revenu espéré 30 et de volatilité 25~;
le trading d'inventaire, 15 et 20~; les solutions structurées, 10 et 12~; corrélations 0.3, 0.2 et 0.4. Son revenu espéré est de 55, sa volatilité de 42.2
et son ratio de Sharpe de 1.30. Des coûts directs de 18 et des \omterm{def:fm:the-head-of-desks-job:budget}{coûts alloués} de 10 laissent, après une rémunération variable de 20~\% d'un résultat positif, un
profit espéré de 20.5, et une probabilité de 26~\% que le revenu ne couvre pas les coûts.

\begin{proposition}[Ce que coûte un objectif au-dessus de l'espérance]\label{prop:fm:the-head-of-desks-job:target}
Soit un revenu annuel suivant une loi normale de moyenne $\mu$ et de volatilité $\sigma$, de ratio de Sharpe $S=\mu/\sigma$, et soit le budget $k\mu$. La probabilité
de l'atteindre vaut
\[
P=\Phi\bigl(S\,(1-k)\bigr).
\]
Pour $k>1$, elle est inférieure à un demi et décroît lorsque le ratio de Sharpe augmente. Multiplier chaque position par $\lambda$ (la moyenne et la volatilité sont toutes deux multipliées) atteint le
budget avec une probabilité $q$ si et seulement si $\lambda\,(\mu-z_q\sigma)=k\mu$, avec $z_q=\Phi^{-1}(q)$, ce qui exige $S>z_q$.
\end{proposition}

\begin{proof}
$P(\text{revenu}\ge k\mu)=1-\Phi((k\mu-\mu)/\sigma)=\Phi(S(1-k))$, décroissante en $S$ lorsque $1-k<0$. Après la mise à l'échelle, le revenu suit une loi normale de moyenne
$\lambda\mu$ et de volatilité $\lambda\sigma$~: $P=q$ équivaut à $(\lambda\mu-k\mu)/(\lambda\sigma)=z_q$, soit $\lambda(\mu-z_q\sigma)=k\mu$, qui admet une
solution positive seulement si $\mu>z_q\sigma$.
\end{proof}

La moitié contre-intuitive de la proposition est la plus utile~: un desk plus régulier a \emph{moins} de chances de battre un objectif ambitieux, car ses
résultats se regroupent plus près de leur moyenne (\cref{fig:fm:the-head-of-desks-job:pmeet}). Pour $k=1.1$ et $S=1.3$, le desk atteint son budget avec
une probabilité de 44.8~\% (45.0~\% sur 20\,000 années simulées)~; pour l'atteindre avec une probabilité de 75~\%, il doit multiplier chaque position par 2.28. Un budget qui se situe
au-dessus du revenu espéré du desk ne demande pas d'effort~: il demande du risque.

\begin{omfigure}
\begin{tikzpicture}
\begin{axis}[width=10.5cm,height=5.4cm,xlabel={\small budget en multiple du revenu espéré $k$},ylabel={\small probabilité de l'atteindre (\%)},
  xmin=0.5,xmax=1.5,ymin=0,ymax=100,tick label style={font=\footnotesize},grid=major,grid style={gray!20},table/col sep=comma,
  legend cell align=left,legend style={font=\scriptsize,at={(0.5,-0.3)},anchor=north,/tikz/every even column/.append style={column sep=8pt}},legend columns=4]
\addplot[omDef,thick] table[x=k,y=s05]{figdata/desk/06-the-head-of-desks-job/pmeet.csv};
\addplot[omProp,thick,dashed] table[x=k,y=s10]{figdata/desk/06-the-head-of-desks-job/pmeet.csv};
\addplot[omThm,thick,dotted] table[x=k,y=s15]{figdata/desk/06-the-head-of-desks-job/pmeet.csv};
\addplot[gray,thick,dashdotted] table[x=k,y=s20]{figdata/desk/06-the-head-of-desks-job/pmeet.csv};
\draw[black,thin] (axis cs:1,0) -- (axis cs:1,100);
\legend{ratio de Sharpe 0.5,ratio de Sharpe 1.0,ratio de Sharpe 1.5,ratio de Sharpe 2.0}
\end{axis}
\end{tikzpicture}

\omcaption{La probabilité que le revenu annuel d'un desk atteigne un budget fixé à $k$ fois son revenu espéré, $\Phi(S(1-k))$, pour quatre
ratios de Sharpe. Toutes les courbes passent par un demi en $k=1$~; au-dessus, plus le desk est régulier, moins il a de chances de battre le budget. Données~:
\texttt{fm\_head.p\_meet\_closed}.}\label{fig:fm:the-head-of-desks-job:pmeet}
\end{omfigure}

\begin{method}[Négocier un budget]\label{met:fm:the-head-of-desks-job:budget}
\begin{enumerate}
\item Estimer le revenu espéré et la volatilité de chaque activité à partir de son propre historique, avec les limites actuelles, et non du seul résultat de l'année précédente.
\item Calculer le revenu espéré, la volatilité et le ratio de Sharpe du desk, ainsi que la probabilité d'atteindre le budget proposé
(\cref{prop:fm:the-head-of-desks-job:target}).
\item Si le budget dépasse l'espérance, énoncer l'échelle de risque qu'il exige et demander les limites qui l'accompagnent, ou un budget
égal à l'espérance.
\item Convenir des \omterm{def:fm:the-head-of-desks-job:budget}{coûts alloués} au desk et de la clé qui les répartit~; vérifier que le desk couvre l'ensemble de ses coûts sur l'année qu'il attend.
\end{enumerate}
\end{method}

\section{Appétit pour le risque~: de la perte que la société peut supporter aux limites}

\begin{definition}[Appétit pour le risque, déclaration d'appétit pour le risque]\label{def:fm:the-head-of-desks-job:appetite}
L'\emph{appétit pour le risque}\index{appétit pour le risque} d'une société est le niveau global et les types de risques qu'elle accepte de prendre, dans la limite du maximum qu'elle pourrait
supporter compte tenu de son capital, de sa liquidité et de ses obligations, pour atteindre sa stratégie et son plan d'activité. Sa \emph{déclaration d'appétit
pour le risque}\index{déclaration d'appétit pour le risque} est la forme écrite~: des énoncés qualitatifs et des mesures quantitatives rapportés aux résultats,
au capital, au risque et à la liquidité, que ses limites de risque répartissent ensuite entre les lignes d'activité, les desks et les types de risque.
\end{definition}

Les définitions reprennent les principes du Financial Stability Board de 2013, qui nomment aussi la borne supérieure, la \emph{capacité de risque}, soit le
risque maximal qu'une société peut prendre avant de franchir ses contraintes réglementaires de capital et de liquidité. Un énoncé tel que «~la société accepte une perte d'au plus
20~millions de dollars en une année sur vingt provenant de ce desk~» se traduit en limites en trois étapes (\cref{fig:fm:the-head-of-desks-job:cascade}).

\begin{omfigure}
\begin{tikzpicture}[font=\footnotesize,>=stealth,box/.style={draw,rounded corners,minimum height=0.7cm,align=center,fill=omDef!10,
  minimum width=3.2cm}]
\node[box,fill=omProp!12] (lt) at (0,2.6) {tolérance aux pertes\\{\scriptsize 20~m\$ sur une année sur vingt}};
\node[box] (vb) at (0,1.2) {budget de volatilité\\{\scriptsize $(\mu+L)/z_{0.95}=45.6$}};
\node[box] (var) at (-3.6,-0.4) {limite de VaR quotidienne\\{\scriptsize $z_{0.99}\,\sigma/\sqrt{252}=6.7$}};
\node[box] (st1) at (0.2,-0.4) {stop~: flux client\\{\scriptsize 8.8}};
\node[box] (st2) at (3.8,-0.4) {stops~: inventaire 7.0,\\{\scriptsize structurées 4.2}};
\draw[->] (lt) -- (vb); \draw[->] (vb) -- (var); \draw[->] (vb) -- (st1); \draw[->] (vb) -- (st2);
\node[font=\scriptsize,align=left,anchor=west] at (2.1,1.2) {volatilité du desk aujourd'hui 42.2~:\\il peut monter de 1.08};
\end{tikzpicture}

\omcaption{La cascade d'\omterm{def:fm:the-head-of-desks-job:appetite}{appétit pour le risque} du desk du chapitre~: une tolérance aux pertes devient un budget de volatilité, qui devient une limite de valeur en risque quotidienne
et un stop par activité, partagés selon la volatilité. En millions de dollars~; desk illustratif. Modèle~: \texttt{firm.deskplan.cascade}.}\label{fig:fm:the-head-of-desks-job:cascade}
\end{omfigure}

Avec un revenu espéré $\mu=55$, une perte d'au plus $L=20$ à 95~\% de confiance autorise une volatilité de $(55+20)/1.645=45.6$~: le desk, à 42.2, peut
monter de 8~\%. À 99~\% sur un jour, la valeur en risque correspondante (livre~6, chapitre~21) vaut $2.326\times45.6/\sqrt{252}=6.7$. Les
20~millions de dollars de tolérance, partagés selon la volatilité de chaque activité, donnent des stops de 8.8, 7.0 et 4.2. La cascade est la hiérarchie de risque du livre~6, chapitre~29, lue
du sommet~: les limites sont la \omterm{def:fm:the-head-of-desks-job:appetite}{déclaration d'appétit pour le risque}, traduite en chiffres.

\omcode{code/firm/deskplan/firm_deskplan.py}{73}{103}{L'échelle de risque qu'exige un budget, et la cascade depuis une tolérance aux pertes jusqu'à un budget de volatilité,
une limite de VaR quotidienne et des stops.}\label{lst:fm:the-head-of-desks-job:cascade}

\begin{remark}[Le budget et l'appétit doivent concorder]\label{rem:fm:the-head-of-desks-job:agree}
L'\omterm{def:fm:the-head-of-desks-job:appetite}{appétit pour le risque} du desk autorise une volatilité de 45.6~; son budget, à tenir trois années sur quatre, exige 2.28 fois 42.2, soit 96. Un plan qui fixe
l'un sans l'autre oblige le responsable de desk à en briser un. Le cycle de planification est le moment où la banque tranche~: un budget plus bas, un appétit plus large
ou un meilleur ratio de Sharpe, et seuls les deux premiers peuvent être retenus en décembre.
\end{remark}

\section{Périmètre de produits et approbation des nouveaux produits}

\begin{definition}[Périmètre de produits, approbation des nouveaux produits]\label{def:fm:the-head-of-desks-job:scope}
Le \emph{périmètre de produits}\index{périmètre de produits} d'un desk est la liste des instruments, des marchés et des types de clients qu'il est autorisé à traiter. L'\emph{approbation des nouveaux produits}\index{approbation des nouveaux produits}
est le processus par lequel une banque examine un produit, un marché ou une activité nouvelle ou modifiée avant qu'un desk puisse
le traiter~: ses risques, sa tarification et sa valorisation, son enregistrement (booking) et son règlement-livraison, son traitement juridique et fiscal, son capital et sa conformité, chacun validé par la
fonction responsable.
\end{definition}

Un nouveau produit est le lieu où les risques d'un desk sont les moins bien compris et ses contrôles les moins éprouvés~: le modèle de prix est neuf, l'enregistrement peut être manuel, le
règlement-livraison n'existe peut-être pas encore. Les superviseurs bancaires américains attendent des banques qu'elles gèrent les risques des «~nouveaux produits, modifiés ou élargis
et services~» avec des approbations appropriées avant qu'une nouvelle activité ne démarre et des contrôles permettant d'identifier, de mesurer, de suivre et de rapporter ses risques. Pour le
responsable de desk, le processus d'approbation est le prix de la croissance~: chaque fonction qui signe est une fonction qui soutiendra le produit le jour où il tournera mal.
Le \omterm{def:fm:the-head-of-desks-job:scope}{périmètre de produits} fixe aussi ce que le desk ne doit \emph{pas} faire~; un trader qui trouve une transaction rentable en dehors de ce périmètre a trouvé une demande d'approbation, et non une
transaction.

\begin{method}[Faire passer un nouveau produit par l'approbation]\label{met:fm:the-head-of-desks-job:npa}
\begin{enumerate}
\item Rédiger la description du produit~: payoff, clients, volume attendu, risque du desk et couvertures.
\item Pour chaque fonction (risque, product control, opérations, juridique, fiscalité, conformité, technologie), lister ce qu'elle doit construire ou vérifier, et obtenir son
accord assorti de conditions.
\item Commencer par des limites de taille et une date de revue~; ne faire passer le produit en activité courante qu'après la première revue.
\end{enumerate}
\end{method}

\section{Le rythme~: le jour, la semaine, l'année}

Un responsable de desk suit un calendrier qui se répète. Le contenu varie d'une banque à l'autre~; la structure est commune, car elle suit les contrôles~:

\begin{center}
\small
\begin{tabular}{@{}lp{10cm}@{}}
\toprule
quand & ce que fait le responsable de desk\\
\midrule
chaque matin & lire le P\&L de la veille et son explication (livre~1, chapitre~7), l'usage des limites et les dépassements~; la réunion du matin sur
les marchés, les positions et les flux clients\\
pendant la journée & approuver les transactions au-delà des limites déléguées~; suivre le risque intrajournalier et les stops~; appels clients\\
chaque soir & valider le P\&L du jour avec le product control (chapitre~15)~; vérifier les positions par rapport aux limites\\
chaque semaine & revoir chaque activité par rapport au plan, à son risque et à son attribution~; recrutement et projets\\
chaque mois & comparer le plan au réalisé (\cref{fig:fm:the-head-of-desks-job:fan})~; la revue du comité des risques (chapitre~12)\\
chaque année & budget et limites~; enveloppe de rémunération et son allocation (chapitre~10)~; \omterm{def:fm:the-head-of-desks-job:scope}{périmètre de produits}\\
\bottomrule
\end{tabular}
\end{center}

La revue mensuelle compare le revenu cumulé au plan. Un plan réparti uniformément sur l'année est une droite~; le revenu du desk est une trajectoire aléatoire
autour de son espérance, et la bande dans laquelle il peut osciller sans signifier quoi que ce soit s'élargit avec la racine carrée du temps
(\cref{fig:fm:the-head-of-desks-job:fan}). Au troisième mois, un desk exactement sur sa trajectoire espérée accuse déjà 1.4~million de dollars de retard sur un budget supérieur de dix
pour cent, un écart largement dans le bruit~; un responsable de desk qui réagit à cet écart en ajoutant du risque réagit au bruit.

\begin{omfigure}
\begin{tikzpicture}
\begin{axis}[width=10.5cm,height=5.6cm,xlabel={\small mois},ylabel={\small revenu cumulé (millions de \$)},xmin=0,xmax=12,ymin=-60,ymax=140,
  xtick={0,2,4,6,8,10,12},tick label style={font=\footnotesize},grid=major,grid style={gray!20},table/col sep=comma,
  legend cell align=left,legend style={font=\scriptsize,at={(0.5,-0.3)},anchor=north,/tikz/every even column/.append style={column sep=8pt}},legend columns=4]
\addplot[name path=a,draw=none,forget plot] table[x=month,y=p5]{figdata/desk/06-the-head-of-desks-job/fan.csv};
\addplot[name path=b,draw=none,forget plot] table[x=month,y=p95]{figdata/desk/06-the-head-of-desks-job/fan.csv};
\addplot[omDef!15] fill between[of=a and b];
\addplot[name path=c,draw=none,forget plot] table[x=month,y=p25]{figdata/desk/06-the-head-of-desks-job/fan.csv};
\addplot[name path=d,draw=none,forget plot] table[x=month,y=p75]{figdata/desk/06-the-head-of-desks-job/fan.csv};
\addplot[omDef!40] fill between[of=c and d];
\addplot[omDef,thick] table[x=month,y=p50]{figdata/desk/06-the-head-of-desks-job/fan.csv};
\addplot[omThm,thick,dashed] table[x=month,y=budget]{figdata/desk/06-the-head-of-desks-job/fan.csv};
\legend{5\up{e}~centile--95\up{e}~centile,25\up{e}~centile--75\up{e}~centile,médiane,budget}
\end{axis}
\end{tikzpicture}

\omcaption{Le desk du chapitre~: centiles du revenu cumulé par mois sur 5\,000 années simulées, face à un budget de 1.1 fois le revenu espéré
revenu réparti uniformément. En décembre, le budget se situe juste au-dessus de la trajectoire médiane. Données~: \texttt{fm\_head.fan}.}\label{fig:fm:the-head-of-desks-job:fan}
\end{omfigure}

\section{Tutoriel~: le plan de l'année prochaine}\label{tut:fm:the-head-of-desks-job:tut}

\begin{tutorial}
\textbf{Objectif.} Construire le plan d'un desk, mesurer la probabilité de l'atteindre et en dériver les limites à partir d'une tolérance aux pertes. \textbf{État final~:}
\cref{fig:fm:the-head-of-desks-job:fan} et les chiffres de la \cref{rem:fm:the-head-of-desks-job:agree}.
\begin{tutsteps}
\item \textbf{Le plan.} \texttt{fm\_head.PLAN} est un \texttt{firm.deskplan.Plan} de trois activités, de leurs corrélations, de leurs coûts et du
budget~; \texttt{plan\_stats} donne un revenu espéré de 55, une volatilité de 42.2, un ratio de Sharpe de 1.30.
\item \textbf{Les cotes.} \texttt{p\_meet} donne 44.8~\%~; \texttt{simulate} confirme le résultat sur 20\,000 années et renvoie les drawdowns maximaux (médiane
de 30.0~millions de dollars, 90\up{e}~centile de 50.4~millions de dollars).
\item \textbf{Le risque que le budget exige.} \texttt{risk\_multiple(PLAN, 0.75)} donne 2.28 (\cref{lst:fm:the-head-of-desks-job:cascade}).
\item \textbf{La cascade.} \texttt{cascade(PLAN, 20)} donne le budget de volatilité, la limite de VaR quotidienne et les stops.
\item \textbf{L'année.} \texttt{fm\_head.fan()} donne les centiles mensuels de la \cref{fig:fm:the-head-of-desks-job:fan}.
\end{tutsteps}
\end{tutorial}

\textbf{Pour aller plus loin.} Porter la corrélation entre flux et inventaire à 0.8 et refaire la cascade~; fixer le budget à l'espérance et
comparer les drawdowns que le desk devra supporter.

\section{Construction~: le plan du desk}\label{bld:fm:the-head-of-desks-job:deskplan}

\begin{build}
\bfield{Rôle} Les chiffres du responsable de desk en un seul endroit~: le plan, ses chances, le risque qu'il exige et les limites que l'appétit de la banque autorise.
\bfield{Interface} \texttt{firm.deskplan}~: \texttt{Business(name, mu, sigma)},
\texttt{Plan(businesses, corr, direct\_cost, allocated\_cost, var\_pay, budget)}~; \texttt{plan\_stats}, \texttt{p\_meet}, \texttt{risk\_multiple(plan, prob)}, \texttt{simulate(plan, n, rng)},
\texttt{cascade(plan, loss\_tolerance, conf, var\_conf)}, \texttt{plan\_vs\_actual(plan, actual\_monthly)}.
\bfield{Règles} Le revenu est normal sur l'année (le chapitre le précise là où c'est utile)~; un budget qui exige un ratio de Sharpe supérieur à $z_q$ renvoie un
multiple de risque infini~; les stops s'additionnent jusqu'à la tolérance aux pertes.
\bfield{Tests d'acceptation} \texttt{code/firm/deskplan/tests/}~: la forme fermée contre la simulation~; le multiple de risque résout son équation et
vaut l'infini quand le ratio de Sharpe est trop faible~; le budget de volatilité et les stops de la cascade~; les lignes plan contre réalisé.
\bfield{Pour aller plus loin} Un revenu à queues épaisses (une loi de Student-$t$) et son effet sur la cascade~; des stops par activité calculés à partir de la contribution de chaque activité à la
volatilité du desk (chapitre~7) plutôt que de sa volatilité autonome.
\end{build}

\begin{omsources}
\item Financial Stability Board, \emph{Principles for an Effective Risk Appetite Framework}, 18 novembre 2013.
\item Office of the Comptroller of the Currency, Bulletin 2017-43, \emph{New, Modified, or Expanded Bank Products and Services: Risk Management
Principles}, 20 octobre 2017.
\end{omsources}

\section{Exercices}

\begin{exercise}[$\star$]\label{exo:fm:the-head-of-desks-job:1}
Un desk de ratio de Sharpe égal à 1 reçoit un budget égal à 1.1 fois son revenu espéré. Calculer sa probabilité de l'atteindre.
\end{exercise}

\begin{exercise}[$\star$]\label{exo:fm:the-head-of-desks-job:2}
Le desk du chapitre attend un revenu de 55 avec une volatilité de 42.2 et supporte des coûts de 28. Calculer la probabilité que le revenu ne couvre pas ses coûts.
\end{exercise}

\begin{exercise}[$\star$]\label{exo:fm:the-head-of-desks-job:3}
Calculer le budget de volatilité d'un desk qui attend un revenu de 55 et dont la firme accepte une perte d'au plus 20 dans une année sur vingt, ainsi que sa limite de VaR journalière à 99~\%.
\end{exercise}

\begin{exercise}[$\star\star$]\label{exo:fm:the-head-of-desks-job:4}
Montrer qu'un desk de ratio de Sharpe 1.3 ne peut atteindre aucun budget avec une probabilité de 95~\% en modifiant la taille de ses positions, quel que soit le budget.
\end{exercise}

\begin{exercise}[$\star\star$]\label{exo:fm:the-head-of-desks-job:5}
À partir de la \cref{fig:fm:the-head-of-desks-job:pmeet}, pour quel multiple de budget un desk de ratio de Sharpe 2 a-t-il la même probabilité de succès qu'un desk
de ratio de Sharpe 0.5 avec $k=1.5$~?
\end{exercise}

\begin{exercise}[$\star\star$]\label{exo:fm:the-head-of-desks-job:6}
Au troisième mois, le desk de la \cref{fig:fm:the-head-of-desks-job:fan} se situe exactement sur sa trajectoire espérée. De combien est-il en retard sur le budget, et à combien
d'écarts types de son revenu sur trois mois cela correspond-il~?
\end{exercise}

\begin{exercise}[$\star\star\star$]\label{exo:fm:the-head-of-desks-job:7}
\emph{Programmation.} Relancer \texttt{firm.deskplan.risk\_multiple} pour le desk du chapitre avec un budget égal à 1.2 fois l'espérance et une
probabilité cible de 75~\%. Quelle échelle est nécessaire, et quelle VaR journalière en découlerait-il~?
\end{exercise}

\begin{exercise}[$\star\star\star$]\label{exo:fm:the-head-of-desks-job:8}
\emph{Trouver la faille.} «~Nous avons manqué le budget deux années de suite alors que notre ratio de Sharpe est de 1.3. Le desk sous-performe.~»
\end{exercise}

\section{Problème~: plus dix pour cent}

\begin{problem}[{Problème du week-end --- plus dix pour cent}]\label{pb:fm:the-head-of-desks-job:1}
Une responsable de desk reçoit le budget de l'année suivante, dix pour cent au-dessus du résultat de cette année, avec des limites inchangées. Elle dispose d'une semaine pour répondre.

\medskip\noindent\textbf{Partie I --- Le plan.}
\begin{enumerate}
\item Définir un \omterm{def:fm:the-head-of-desks-job:budget}{budget de revenus} et un \omterm{def:fm:the-head-of-desks-job:budget}{coût alloué}.
\item Donner le revenu espéré, la volatilité et le ratio de Sharpe du desk à partir de ses trois activités.
\item Donner son profit espéré après coûts et rémunération variable, ainsi que la probabilité que le revenu ne couvre pas les coûts.
\item Quel est le budget, en millions de \$~?
\end{enumerate}

\medskip\noindent\textbf{Partie II --- Les cotes.}
\begin{enumerate}[resume]
\item Énoncer et démontrer la \cref{prop:fm:the-head-of-desks-job:target}.
\item Donner la probabilité d'atteindre le budget, sous forme fermée et par simulation.
\item Pourquoi un desk plus régulier a-t-il moins de chances d'atteindre un budget ambitieux~?
\item Donner la médiane et le 90\up{e}~centile du drawdown maximal du desk sur une année.
\item Quelle échelle de risque permet d'atteindre le budget trois années sur quatre~?
\end{enumerate}

\medskip\noindent\textbf{Partie III --- L'appétit.}
\begin{enumerate}[resume]
\item Définir l'\omterm{def:fm:the-head-of-desks-job:appetite}{appétit pour le risque} et la \omterm{def:fm:the-head-of-desks-job:appetite}{déclaration d'appétit pour le risque}, et nommer la borne supérieure de cet appétit.
\item Convertir une tolérance de perte de 20~millions de \$ à 95~\% en budget de volatilité.
\item Donner la limite de VaR journalière à 99~\% et les stops par activité.
\item De combien le desk peut-il augmenter son échelle dans le cadre de son appétit~?
\item Comparer le risque que permet l'appétit avec le risque qu'exige le budget.
\end{enumerate}

\medskip\noindent\textbf{Partie IV --- La réponse.}
\begin{enumerate}[resume]
\item Définir le \omterm{def:fm:the-head-of-desks-job:scope}{périmètre de produits} et l'\omterm{def:fm:the-head-of-desks-job:scope}{approbation des nouveaux produits}, et expliquer pourquoi un nouveau produit est l'endroit où les contrôles sont les plus faibles.
\item Au troisième mois, le desk se situe sur sa trajectoire espérée. De combien est-il en retard sur le budget, en millions de \$ et en écarts types~?
\item Que la responsable de desk ne devrait-elle pas faire en réponse~?
\item Lesquelles des entrées du plan peuvent être modifiées en décembre, et lesquelles ne le peuvent pas~?
\item Énoncer le \emph{résultat nommé}~: la probabilité d'atteindre un budget de $k$ fois le revenu espéré au ratio de Sharpe $S$, sa valeur pour ce
desk, et l'échelle de risque nécessaire pour l'atteindre trois années sur quatre.
\item En deux phrases, rédiger sa réponse au budget.
\end{enumerate}
\end{problem}

\section{Questions d'entretien}

\begin{interviewq}[$\star$ \iqroles{trader}]\label{iq:fm:the-head-of-desks-job:1}
Le ratio de Sharpe de votre desk est de 1. Quelle est la probabilité qu'il gagne moins que son revenu espéré sur une année~? Moins que zéro~?
\end{interviewq}

\begin{interviewq}[$\star$ \iqroles{risque}]\label{iq:fm:the-head-of-desks-job:2}
Qu'est-ce qu'une \omterm{def:fm:the-head-of-desks-job:appetite}{déclaration d'appétit pour le risque}, et comment devient-elle les limites d'un desk~?
\end{interviewq}

\begin{interviewq}[$\star\star$ \iqroles{trader, chercheur}]\label{iq:fm:the-head-of-desks-job:3}
La direction augmente votre budget de 20~\% et maintient vos limites. Que fait cela à vos cotes, et que demanderiez-vous~?
\end{interviewq}

\begin{interviewq}[$\star\star$ \iqroles{trader}]\label{iq:fm:the-head-of-desks-job:4}
Vous êtes au troisième mois de l'année et en retard sur un budget réparti uniformément. Quand cela devrait-il vous inquiéter~?
\end{interviewq}

\begin{interviewq}[$\star\star$ \iqroles{banque, risque}]\label{iq:fm:the-head-of-desks-job:5}
Un trader souhaite négocier un produit que le desk n'a jamais traité. Que doit-il se passer en premier, et pourquoi~?
\end{interviewq}

\begin{interviewq}[$\star\star\star$ \iqroles{chercheur, risque}]\label{iq:fm:the-head-of-desks-job:6}
Pourquoi un ratio de Sharpe plus élevé peut-il réduire la probabilité d'atteindre un budget~? Quand l'augmente-t-il~?
\end{interviewq}
